\documentclass[10pt,a4paper]{amsart}
\usepackage[margin=19mm]{geometry}
\usepackage{amsmath,amssymb,mathtools}
\usepackage[scr=rsfs,cal=euler]{mathalfa}
\usepackage{xfrac}
\usepackage[unicode,hidelinks,bookmarks=false]{hyperref}
\newcommand{\ie}{i.e.\ }
\newcommand{\cf}{cf.\ }
\newcommand{\resp}{respectively\ }
\newcommand{\Subsection}{\subsection}
\providecommand{\nobreakdash}{\nobreak\hbox{-}}
\setlength{\emergencystretch}{2em}
\multlinegap0pt
\usepackage{bm}
\usepackage{bbm}
\usepackage{dsfont}
\usepackage{xspace}
\usepackage{cancel}
\usepackage{mathtools}
\usepackage{amsthm}
\makeatletter
\@ifundefined{Theorem}{\newtheorem{Theorem}{Theorem}}{}
\@ifundefined{Lemma}{\newtheorem{Lemma}[Theorem]{Lemma}}{}
\makeatother
\newcommand{\R}{\mathbb R}
\newcommand{\Z}{\mathbb Z}
\newcommand{\conv}{\operatorname{conv}}
\newcommand{\ls}{\operatorname{ls_\Delta}}
\newcommand{\w}{\operatorname{w}}
\title[Lattice Size in Higher Dimension]{Lattice Size in Higher Dimension}
\author{Abdulrahman Alajmi, Sayok Chakravarty, Zachary Kaplan, Jenya Soprunova}
\date{Source: arXiv:2209.00712v2 (CC BY 4.0)}
\begin{document}
\maketitle

\noindent\emph{Source and licence.} Lattice Size in Higher Dimension. Version arXiv:2209.00712v2 (CC BY 4.0), distributed under the Creative Commons Attribution 4.0 International licence, as recorded in the arXiv licence metadata for this exact version and shown on the abstract page https://arxiv.org/abs/2209.00712v2. Full rights notice: licenses/arxiv-2209.00712-CC-BY-4.0.txt.\par\smallskip

\noindent\textit{Editorial scope.} Counted unit: the Theorem whose source label is \texttt{T:Main1} in the article (source0.tex lines 450--453, source label \texttt{T:Main1}), delivered together with the article's own complete original proof of it (source0.tex lines 455--510, one \texttt{proof} environment of 56 lines). No mathematics is written, repaired or re-derived: statement and proof are the article's own text line for line. Required context retained uncounted: L:IntroLemma1 (lines 297--306); L:IntroLemma2 (lines 331--343); L:Lemma1 (lines 408--410); L:Lemma2 (lines 431--433). Every definition, display and label of those lines is retained unchanged. Omitted from this package and not redistributed: 1--296, 307--330, 344--407, 411--430, 434--449, 454--454, 511--598. Nothing inside a retained excerpt is omitted or altered: every retained body is delivered exactly as the article's lines are written, so no omission receipt of any kind accompanies this package. Citations: the retained excerpts carry no citation command at all, so no bibliography entry of the article is transcribed and no reference is redistributed. Reference adaptation: none was needed and none was made; every reference inside the delivered unit resolves inside it, so no unresolved reference or citation is produced. Printed slips kept literal: no printed word, symbol, hypothesis or number of the counted statement or of its proof was altered. Layout and class: the article's own class is \texttt{amsart} with packages including inputenc, amssymb, hyperref, multicol, colordvi, graphicx, pgf, epsfig, epsf; the wrapper is the lane's pinned \texttt{amsart} editorial wrapper at 10pt on A4, which restates the theorem-like environments the retained text needs and adds the article's own macro definitions that the retained text needs (\texttt{\textbackslash R}, \texttt{\textbackslash Z}, \texttt{\textbackslash conv}, \texttt{\textbackslash ls}, \texttt{\textbackslash w}), with extra wrapper packages (bm, bbm, dsfont, xspace, cancel, mathtools). The delivered numbering is the unit's own. Assets: no retained span contains a figure environment, a graphics-inclusion command, a PDF-page inclusion, a table or a TikZ picture; no image file is copied into this package, the package declares no asset dependency and no image file is redistributed with it. Licence: version arXiv:2209.00712v2 of this article is distributed under the Creative Commons Attribution 4.0 International licence, as recorded in the arXiv licence metadata for this exact version and shown on the abstract page https://arxiv.org/abs/2209.00712v2. A full rights notice is delivered at licenses/arxiv-2209.00712-CC-BY-4.0.txt. Characters: every retained excerpt is pure ASCII, so the delivered engine, XeLaTeX with its default Latin Modern text font, reports no missing character.
\begin{Lemma}\label{L:IntroLemma1}
\begin{itemize} 
\item[(1)] For any $h\in \Z^d$ we have $\w_{h}(AP)=\w_{A^{T}h}(P)$.
\item[(2)] For $i=1,\dots, d$ we have $\w_{e_i}(AP)=\w_{h_i}(P)$.
\item[(3)] Let $(e_1,\dots,e_d)$ be the standard basis of $\R^d$. Then $\w_{e_i}(P)\leq l_1(P)$ for $i=1,\dots, d$. 
\item[(4)] For $i=1,\dots, d$ we have $\w_{h_i}(P)\leq l_1(AP)$.
\item[(5)] Let $e\in\Z^d$ be a vector whose entries lie  in $\{0,1\}$. Then $\w_{e}(P)\leq l_1(P)$.
\item[(6)] Let $h$ be the sum of any non-empty collection of rows of $A$. Then $\w_{h}(P)\leq l_1(AP)$.
\end{itemize}
\end{Lemma}

\begin{Lemma}\label{L:IntroLemma2} We have
\begin{itemize} 
\item[(1)] 
$l_1(AP)=\max_{x\in P}\langle h_1+\cdots+h_d,x\rangle-\min_{x\in P}\langle h_1,x\rangle-\cdots-\min_{x\in P}\langle h_d,x\rangle$;
\item[(2)] $l_1(AP)$ does not depend on the order of rows in $A$;
\item[(3)] $
l_1(AP)=l_1(BP)$, where $B=\left(\begin{matrix}h_1\\ \vdots\\h_{d-1}\\
%-(h_1+\cdots+h_{d-1}+h_d)
-\sum_{i=1}^d h_i
\end{matrix}\right).
$
\end{itemize}
\end{Lemma}

\begin{Lemma}\label{L:Lemma1} Let $h=(a_1,\dots, a_{d+1})\in\R^{d+1}$ be a primitive vector with $\w_{h}(T_{p_1\dots p_d})\leq k+2$.
Then $a_d\in\{0,1,-1\}$.
\end{Lemma}

\begin{Lemma}\label{L:Lemma2} Let $h=(a_1,\dots, a_{d+1})\in\R^{d+1}$ be a primitive vector with $a_d=\pm 1$ and  $\w_{h}(T_{p_1\dots p_d})\leq  k+2$.
Then $\w_{h}(T_{p_1\dots p_d})= k+2$.
\end{Lemma}

\begin{Theorem}\label{T:Main1} 
Let $\alpha=p_1+\cdots+p_{d-1}$, where all $p_i$ are positive and $p_d\geq 2$. Define $k=\lfloor\frac{p_d-2}{\alpha+1}\rfloor$.
Suppose that $p_d\geq \alpha^2-\alpha$. Then $\ls(T_{p_1\dots p_d})=k+3$.
\end{Theorem}

\begin{proof}
Suppose that there exists a unimodular map $L$ that maps $T_{p_1\dots p_d}$ inside $(k+2)\Delta$ and let $A$ be the corresponding 
unimodular matrix. Then by Lemma~\ref{L:IntroLemma1} for each of its rows  $h$ we have $\w_{h}(T_{p_1\dots p_d})\leq k+2$. We also have the same 
inequality for the sum of any nonempty collection of rows of $A$. By Lemma~\ref{L:Lemma1} each of the entries in the $d$th column of $A$
is 0, 1, or $-1$, and the same applies to the sum of any collection of entries in the  $d$th column of $A$. Hence, up to permutation of rows, the $d$th column
of $A$ is $(0,0,\dots, 0,\pm 1)^T$ or $(0,0,\dots, 0,1, -1)^T$ and by Lemma~\ref{L:IntroLemma2} we can assume that it is the former. Then by Lemma~\ref{L:Lemma2} we have $\w_{e_{d+1}}(L(T_{p_1\dots p_d}))=k+2$.
Since $L(T_{p_1\dots p_d})\subset (k+2)\Delta$ we conclude that 
$$(0,\dots, k+2)\in L(T_{p_1\dots p_d}).$$ 
Similarly, we have $\w_{e_1+\cdots+e_{d+1}}(L(T_{p_1\dots p_d}))=k+2$ and, together 
with $L(T_{p_1\dots p_d})\subset (k+2)\Delta$, this implies that $L(T_{p_1\dots p_d})$ contains the origin. Hence $L(T_{p_1\dots p_d})$ and, therefore, $T_{p_1\dots p_d}$ contains an edge of lattice length 
$k+2$. All the edges of $T_{p_1\dots p_d}$ are primitive, except, possibly, for the one connecting points $(0,\dots,1)$ and $(p_1,\dots,p_d,1)$, whose lattice length is $\gcd(p_1,\dots,p_d)$.
Hence we conclude that $\gcd(p_1,\dots,p_d)=k+2$.   Using the assumption $p_d\geq \alpha^2-\alpha$ we get
$$k+2=\gcd(p_1,\dots, p_d)\leq p_1+\dots+p_{d-1}=\alpha\leq \frac{p_d-2}{\alpha+1}+2<k+3.
$$

Note that since all the $p_i$ are positive, we have $\gcd(p_1,\dots, p_d)< p_1+\dots+p_{d-1}$ unless $d=2$ and $p_2$ is a multiple of of $p_1$.
When the inequality is strict we arrive  at a contradiction since then integer $p_1+\dots+p_{d-1}$ is strictly between the consecutive integers $k+2$ and $k+3$.

It remains to consider the case when $d=2$, $p_2$ is a multiple of $p_1$, and $k+2=\alpha=p_1$.  
We have $\lfloor\frac{p_2-2}{p_1+1}\rfloor=k=p_1-2$ and hence
$$p_2-2=(p_1+1)(p_1-2)+r,
$$
where $0\leq r\leq p_1$. We get  $p_2=p_1^2-p_1+r$ and, since $p_2$ is a multiple of $p_1$, there are two options: $p_2=p_1^2$ and $p_2=p_1^2-p_1$.

Suppose first $p_2=p_1^2$ and let $h=(a_1, a_2,a_3)$ be a direction with $\w_h(T_{p_1p_2})\leq k+2=p_1$.
By Lemma~\ref{L:Lemma1} we have $a_2=0,\pm 1$.  For  $a_2=-1$ we get
$$\max\{a_1,-1,a_3,a_1p_1-p_1^2+a_3\}-\min\{a_1,-1,a_3,a_1p_1-p_1^2+a_3\}\le p_1,
$$
which implies $a_1+1\leq p_1$ and $-a_1p_1+p_1^2\leq p_1$, so we conclude that $a_1=p_1-1$. Plugging in this value for $a_1$ we get
$$\max\{p_1-1,-1,a_3,a_3-p_1\}-\min\{p_1-1,-1,a_3,a_3-p_1\}\le p_1,
$$
which implies $a_3+1\leq p_1$ and $p_1-1-(a_3-p_1)\leq p_1$, so $a_3=p_1-1$. We conclude that for $a_2=-1$ the only $h$ with $\w_h(T_{p_1p_2})\leq p_1$ is 
$h=(p_1-1,-1,p_1-1)$ and for such $h$ we get $\w_h(T_{p_1p_2})=k+2$. Similarly, for $a_2=1$ such direction is $h=(1-p_1,1,1-p_1)$.



Hence  in this case we can only use as rows of $A$ vectors $\pm (p_1-1,-1,p_1-1)$ and vectors whose second component is 0.
Further, we can assume that the second column of $A$ is $(0,0,\pm 1)^T$, and the third row is $\pm (p_1-1,-1,p_1-1)$.
Then the sum of the third row with any of the first two rows will also have to be of the same form as the third row, which would imply $\det A=0$.

We next consider the last case $p_2=p_1^2-p_1$. Recall that we have $k+2=p_1$. By Lemma~\ref{L:Lemma1} and Lemma~\ref{L:Lemma2} if $\w_h(T_{p_1p_2})\leq p_1$ for 
$h=(a_1,a_2,a_3)$ then $a_2=0,\pm 1$ and if $a_2=\pm 1$ we have $\w_h(T_{p_1p_2})= p_1$. Let's further investigate the case $a_2=0$. We have 
$$\max\{a_1,0,a_3,a_1p_1+a_3\}-\min\{a_1,0,a_3,a_1p_1+a_3\}\le p_1,
$$
where we can assume $a_1\geq 0$. Hence $a_1p_1\leq p_1$, which implies $a_1=0$ or $1$ and in the latter case the width is $p_1$. We conclude that the only direction $h$ with $\w_h(T_{p_1p_2})< p_1$ is 
$(0,0,\pm 1)$.

As before, we can assume that the second column in matrix $A$ is $(0,0,1)^T$. Denote the rows of $A$ by $h_1,h_2$, and $h_3$.
Then out of the widths of $T_{p_1p_2}$ in the direction of  $h_1$ and $h_2$ only one can be strictly less than $p_1$.
We can assume that this happens in the direction of $h_1$. We also know that the width of $T_{p_1p_2}$ in the direction of $h_3$ and $h_1+h_2+h_3$
is $p_1$. We can now conclude that $L(T_{p_1p_2})$ contains the triangle 
$$\conv\{(0,0,0), (0,p_1,0), (0,0,p_1)\}.
$$
Note that we assumed that $p_d\geq 2$, which implies $p_1^2-p_1=p_2\geq 2$ and hence $p_1\geq 2$. Hence our conclusion implies  that  $T_{p_1p_2}$ contains a lattice triangle with three 
non-primitive sides, and this contradiction completes the argument.
\end{proof}

\end{document}
